\chapter*{Lời Mở Đầu}
\addcontentsline{toc}{chapter}{Lời Mở Đầu}
\markboth{Lời Mở Đầu}{Lời Mở Đầu}

\begin{quotecard}{Một cuốn sách đọc xong đừng chỉ cười}
Tên sách có thể bắt trend. Nhưng thứ giữ người đọc ở lại không phải cái tên. Đó là cảm giác: “Ủa, câu này mình từng nói thật.” Rồi thêm một nhịp nữa: “Ừ, nghe đáp lại vậy thì cũng khó cãi tiếp.”

\textit{(A trendy title may pull attention first. But what keeps readers is the moment they realize: “I have actually said something like this.” Then comes the second beat: “Fair enough. That reply leaves me little room to argue.”)}
\end{quotecard}

\begin{truyen}{Cuốn này sinh ra để dùng rất thật}{Why this edition exists}
\chuhoa{C}{uốn này không được thiết kế để ngồi trên kệ cho đẹp. Nó được thiết kế để mở ra giữa lớp học, giữa giờ sinh hoạt, trước một bài kiểm tra, hoặc đúng lúc học sinh đang chuẩn bị thả ra một câu chống chế nghe rất mượt mà.}
\textit{(This edition is not built to sit prettily on a shelf. It is built to be opened in class, in homeroom, before a test, or right when a student is about to launch a beautifully polished excuse.)}

Điểm mới của bản này là từng câu hỏi - đáp được tách riêng thành từng box. Nghĩa là mỗi pha đối thoại đều tự đứng được như một “màn”. Giáo viên có thể bốc ngẫu nhiên một câu để mở tiết. Học sinh có thể đọc đúng một câu rồi tự ái nhẹ một chút, nhưng nhờ vậy mà nhớ lâu hơn.
\textit{(In this edition, every question and reply stands in its own visual box. Each exchange becomes a small self-contained scene. Teachers can pick one at random to open a lesson. Students can read one line, feel a little exposed, and therefore remember it longer.)}

Giọng sách vẫn giữ đúng nguyên tắc cũ: khịa thật, vui miệng, nhưng không rẻ. Đáp mạnh, nhưng để kéo người đọc tỉnh ra, chứ không phải để làm họ bé đi.
\textit{(The voice keeps the same rule: sharp, playful, but not cheap. Strong replies are here to wake readers up, not to reduce them.)}
\end{truyen}

\begin{baihoc}
Một câu đáp hay không phải câu làm người nghe quê nhất. Nó là câu làm người nghe im lại lâu nhất vì thấy đúng.

\textit{(The best comeback is not the one that embarrasses the listener most. It is the one that leaves them silent the longest because it is true.)}
\end{baihoc}

\begin{ghinhoanh}
\textbf{Short English to remember:}

Sharp is useful only when it reveals truth.

If it only stings, it is not deep enough yet.
\end{ghinhoanh}

\begin{tuhocnhanh}
1. Em hay chống chế vì lười, vì sợ, hay vì muốn giữ mặt? \textit{(Do you make excuses out of laziness, fear, or the need to save face?)}

2. Một câu em từng nói tưởng rất khôn nhưng giờ nghĩ lại hơi quê là gì? \textit{(What is one sentence you once thought was clever, but now feels a little embarrassing?)}

3. Nếu bị bóc đúng chỗ đau, em muốn sửa thật hay cãi tiếp cho đỡ ngượng? \textit{(When someone hits the painful truth exactly, do you want to improve or keep arguing to avoid embarrassment?)}
\end{tuhocnhanh}

\ngancach